\documentclass[12pt]{article}
\usepackage[utf8]{inputenc}
\usepackage[T1]{fontenc}
\usepackage{hyperref}
\title{Tuvaluan Language Resources for K-12 Curriculum Development}
\author{Compiled by Claude Sonnet 5 (claude-sonnet-5) for Andr\'as Kornai}
\date{}
\begin{document}
\maketitle

\section{Introduction}
Tuvaluan (ISO 639-3 tvl; Glottocode tuva1244) is a Polynesian language of the Austronesian family, spoken by about 13,000 people. It is used as a language of instruction in Tuvalu. This report is based on a small number of web searches, and little curriculum material turned up.

\section{Primary Educational Resources}
A learner.net page says Tuvaluan is the main language of instruction in the first three primary years, with English taking over from Year 4. The site is of uncertain reliability, so check the claim against Asian Development Bank or UNESCO documents. The page says a Class 1 Tuvaluan syllabus was not readily found online.
\begin{itemize}
\item Tuvaluan Language Learning and Education in Tuvalu: \url{https://learner.net/u/tuvalu-grade-1-tuvaluan-language}
\item Tuvalu Primary School Curriculum - English: \url{https://learner.net/u/tuvalu-year-three-primary-english}
\end{itemize}

\section{Linguistic and Reference Materials}
\begin{itemize}
\item Dictionary of the Tuvaluan Language and Culture (Tachibana and Hoshi, ILCAA, online, English and Japanese, still being updated): \url{https://tuvalu.aa-ken.jp/en/}
\item Reference list for that dictionary, which includes Besnier's \emph{Tuvaluan: A Polynesian Language of the Central Pacific} (Routledge, 2002), Jackson's \emph{Tuvaluan Dictionary} and an unpublished primary-school textbook by Teagai Esekia: \url{https://tuvalu.aa-ken.jp/en/reference}
\item A Tuvaluan Lexicon of the Coconut Culture: \url{https://irc.aa.tufs.ac.jp/en/announce/topics/2021-07-20-1115/}
\item Entries from Jackson 2001 in Pollex-Online: \url{https://pollex.eva.mpg.de/source/35/?page=5}
\end{itemize}

\section{Religious and Community Resources}
No religious or community resources were found in these searches.

\section{Digital and Audio Resources}
The online dictionary above includes photos and illustrations. Related ILCAA tools and archives are listed at \url{https://lingdy.aa-ken.jp/en/publications/200208-tools-and-archives/}. No audio resources were confirmed.

\section{Curriculum Development Context}
Tuvaluan is already the language of instruction in the early primary years. The main gaps found are published Tuvaluan textbooks for later years and an accessible official syllabus. Next steps are to check the Tuvalu Ministry of Education, ADB and UNESCO documents, and local publishers.
\end{document}